\documentclass[11pt]{article}
\usepackage{kristiania-workbook}

\begin{document}

\begin{center}
  {\LARGE\bfseries Profile your Linux system -- the baseline}\\[3pt]
  {\large Session 2 -- Linux \& cloud infrastructure foundations}\\[5pt]
  {\small Exercise 2 \textperiodcentered\ Session 2 lab \textperiodcentered\ Estimated time: 60--90 min}
\end{center}
\vspace{2pt}\hrule\vspace{1.1em}

Write your answers in the answer document \cmd{SKY2100\_Exercise2\_Answers.docx}. You cannot detect something \emph{abnormal} until you
know what \emph{normal} looks like. Today you measure your VM across the four pillars (CPU, memory,
disk, network) plus its logs, and record a \textbf{baseline} you will reuse all semester. These are
the same instruments that later feed your IDS (Session~9) and your incident timeline (Session~10).

\wbsection{1\quad Measure the four pillars}

\task{Install the monitors:} \cmd{sudo apt update \&\& sudo apt install -y htop btop}
\note{If \cmd{btop} is unavailable on your release, use \cmd{htop}; if neither installs, \cmd{top} is
always present. Record which one you used in the answer document.}

\task{Run each command, read the output, and record the reading in the baseline table of the answer
document (same five rows):}
\begin{itemize}
\item CPU model \& cores --- \cmd{lscpu}
\item Total / available RAM --- \cmd{free -h}
\item Disk total / free (\cmd{/}) --- \cmd{df -h}
\item Listening ports (count + which) --- \cmd{ss -tulpn}
\item Current load average --- \cmd{top} (or \cmd{htop} / \cmd{btop})
\end{itemize}

\task{Read the recent log as well:} \cmd{journalctl -xe | tail -n 30} --- nothing to record, but you will
need this command again in Sessions 9 and 10.

\wbsection{2\quad Cloud transfer: Microsoft Learn}
\task{Complete ``Describe the core architectural components of Azure''} (Part~2 of the series).
\par\noindent\mbox{\href{https://learn.microsoft.com/training/modules/describe-core-architectural-components-of-azure/}{learn.microsoft.com/training/modules/describe-core-architectural-components-of-azure}}

\task{Map the concepts.} In the answer document, write one line each connecting the training to the
lecture: Region, Availability Zone, Resource group, Subscription, Azure Resource Manager (ARM).

\wbsection{3\quad Azure reflection}
\task{Your VM has a fixed CPU/RAM/disk. How does an Azure \emph{region} and \emph{availability zone}
change the availability picture compared with your single local machine?}

\task{\cmd{ss -tulpn} shows the ports your VM exposes. In Azure, which construct would limit who can
reach those ports, and at which level (host vs.\ network)?}

\aha{The numbers you just wrote down are not interesting on their own. They matter because they are
what ``normal'' looks like on your machine, and you can only recognise trouble by comparing against
them. When Session~9 has your intrusion detection flag a process, or Session~10 asks you to explain a
spike in a log, the first question is always the same: is this different from the baseline? An unknown
port in \cmd{ss}, a load average far above today's, or free memory near zero all mean something only
because you measured the quiet state first. Keep this baseline; you will compare against it for the rest
of the course.}

\wbsection{4\quad Knowledge check}
\begin{enumerate}
\item Which command lists the ports your machine is \emph{listening} on?\\
      \cbox \cmd{free -h} \quad \cbox \cmd{ss -tulpn} \quad \cbox \cmd{df -h} \quad \cbox \cmd{lscpu}
\hint{The socket-statistics command you ran for the baseline today.}
\item State the difference between a \textbf{Type 1} and a \textbf{Type 2} hypervisor, with one
      example of each.
\hint{One runs directly on the hardware, one runs as an app on a host OS.}
\item In a data centre, \textbf{east-west} traffic means\dots\\
      \cbox client $\leftrightarrow$ data centre \quad
      \cbox server $\leftrightarrow$ server inside the data centre \quad
      \cbox data centre $\leftrightarrow$ internet
\hint{Most data-centre traffic stays server-to-server, inside the building.}
\item Comer notes that running several VMs on one server adds overhead. Give \emph{two} sources of
      that overhead.
\hint{Each VM carries its own OS — its own scheduler, background processes, memory.}
\item An \textbf{Availability Zone} is best described as\dots\\
      \cbox a whole region \quad \cbox a physically separate data centre within a region \quad
      \cbox a resource group \quad \cbox a subscription
\hint{A physically separate data centre inside one region.}
\item Put the Azure hierarchy in order from \emph{broadest} to \emph{narrowest}: resource, management
      group, subscription, resource group.
\hint{From the top down: management group, then narrower, to the single resource.}
\item Why is a virtual machine described as ``a digital object'', and what is \emph{one} security risk
      that follows from it?
\hint{Its entire state is data — which means it can be saved to a file, and copied.}
\item In \cmd{free -h}, which value best indicates memory \emph{pressure}?\\
      \cbox total \quad \cbox available / swap in use \quad \cbox buffers \quad \cbox shared
\hint{Not the 'free' column — the one that counts cache as usable.}
\item Explain \textbf{hot-aisle / cold-aisle} containment in one sentence, and say which CIA property
      it ultimately protects.
\hint{It is about airflow: keep cold intake and hot exhaust apart — which property does overheating threaten?}
\item What is a \textbf{hypervisor}, and which CPU feature lets it stay more privileged than the guest
      operating systems?
\hint{It creates and mediates the VMs; CPU rings let it outrank the guests.}
\item Name the three persistent-storage abstractions used in cloud data centres (Comer Ch.~8).
\hint{Block, file, and the key-value one.}
\item (Reflection) Which single base-metric command will you rely on most when hunting an incident,
      and why?
\end{enumerate}

\signoff

\clearpage
\solheader
{\LARGE\bfseries\color{kred} Solutions}\\[2pt]
{\small Work through the lab first, then check here. Model answers are indicative — equivalent reasoning is fine.}\par\vspace{8pt}
\noindent\textbf{Note:} Model answers are indicative -- accept equivalent reasoning. Multiple-choice
answers are in \textbf{bold}.

\wbsection{Knowledge check}
\begin{enumerate}
\item \textbf{\cmd{ss -tulpn}}.
\item \textbf{Type 1} (bare-metal) runs directly on hardware -- ESXi, Xen, Hyper-V; \textbf{Type 2}
      (hosted) runs as an app on a host OS -- VMware Workstation, VirtualBox.
\item \textbf{Server $\leftrightarrow$ server inside the data centre.}
\item Each VM runs its \emph{own} OS, so each adds a \emph{scheduler} and a set of \emph{background
      processes}; plus per-VM memory/CPU and the time to boot an OS. (Comer Ch.~6, p.~71--72,
      verified.)
\item \textbf{A physically separate data centre within a region.}
\item \textbf{Management group $\rightarrow$ subscription $\rightarrow$ resource group $\rightarrow$
      resource.}
\item A VM's entire state (memory, disk, devices) is data, so it can be saved as a file/snapshot;
      risk: it can be \emph{copied or stolen wholesale}, and a snapshot may capture secrets in memory.
      (Comer Ch.~5, p.~63, verified.)
\item \textbf{Available memory / swap in use.}
\item Racks are arranged so \emph{cold intake} and \emph{hot exhaust} air never mix, improving cooling
      efficiency; it ultimately protects \textbf{Availability} (keeps hardware running). (Comer p.~40,
      verified.)
\item A \textbf{hypervisor} creates and mediates VMs; CPU \textbf{privilege levels (rings)} let it run
      more privileged than the demoted guest OSes. (Comer p.~59--60, verified.)
\item \textbf{Block, file, and object (key-value) storage.} (Comer Ch.~8, verified.)
\item Reflection -- commonly \cmd{journalctl} (the log source for the timeline) or \cmd{ss} (attack
      surface).
\end{enumerate}

\wbsection{References}
Comer Ch.~4 (hot/cold aisles \textbf{p.~40}, east-west traffic \textbf{p.~44}, leaf-spine
\textbf{p.~47}) \& Ch.~5 (approaches \textbf{p.~55}, privilege levels \textbf{p.~59}, extending
privilege \textbf{p.~60}, levels of trust \textbf{p.~60}, VM as digital object \textbf{p.~63}, VM
migration \textbf{p.~64}, live migration \textbf{p.~65}); Ch.~6 VM overhead \textbf{p.~71}; Ch.~8
storage abstractions -- all verified. J\o sang Ch.~3 (syllabus, page numbers not verified). MS Learn:
\emph{Describe the core architectural components of Azure} (verified, June 2026).

\end{document}
